% hci-spine LaTeX skeleton — contribution type: EMPIRICAL (user study)
% Venues: CHI, CSCW. Wobbrock class: Empirical.
\documentclass[sigconf,review,anonymous]{acmart}
\settopmatter{printacmref=false}
\usepackage{booktabs}
\begin{document}

\title{TITLE — name the phenomenon / question}
\begin{abstract}
% Phenomenon + RQs + what was found. Universal check: N here == N in Method == N in tables.
\end{abstract}
\maketitle

\section{Introduction}
% Phenomenon + explicit Research Questions (RQ1 ...). End with contribution list.

\section{Related Work}

\section{Method}
% STRICT honesty zone. Study design matching the RQ; measures' validity.
\subsection{Participants}
% Recruitment, demographics, justification of N. Bias acknowledged. Ethics/IRB.
\subsection{Procedure}
\subsection{Measures}
\subsection{Analysis plan}
% Correct test choice / qualitative coding process + reliability (IRR or reflexive-TA stance).
% Universal check: no HARKing — RQs/hypotheses a priori or disclosed exploratory.

\section{Findings}
% Organized by RQ or theme. Honesty: every thematic claim names its evidence
% (which/how many participants); quotes verbatim + attributed to real IDs;
% every statistic recomputable (exact test, df, p, effect size); report nulls.
\begin{table}[t]\centering\caption{Participants / conditions.}
\begin{tabular}{lll}\toprule & & \\ \bottomrule\end{tabular}\end{table}

\section{Discussion}
% Interpretation + design implications. Stay within what the data supports.
\section{Limitations}
\section{Conclusion}

\bibliographystyle{ACM-Reference-Format}
\bibliography{refs}
\end{document}
